\subsection{Construction Mechanisms (RQ2)}
\label{sec:ablations}

We examine support learning, distribution matching, feature construction, and assignment sharing on SportsFit and BookHis, using one class order, three construction seeds, and three student seeds per memory.

\begin{table}[t]
\centering
\caption{Core construction controls. Each entry reports Macro-F1 / Accuracy.}
\label{tab:loss_controls}
\begin{tabular}{ccc}
\toprule
Constructor & SportsFit & BookHis\\
\midrule
\method{} & 80.02 / 85.38 & 48.53 / 81.38\\
Fixed \(Q\) & 58.35 / 67.23 & 39.62 / 73.97\\
Without DM & 65.28 / 72.27 & 47.23 / 79.43\\
Without GM & 80.47 / 85.81 & 48.47 / 81.31\\
\bottomrule
\end{tabular}
\end{table}

Keeping \(Q\) fixed lowers F1 by 21.67 points on SportsFit and 8.91 on BookHis. Removing distribution matching (DM) lowers it by 14.74 and 1.30 points. These controls identify learned source contributions and class matching as the clearest gains. Removing gradient matching (GM) produces similar scores, so GM is an auxiliary feedback path rather than a necessary condition for predictive performance in this batch.

\noindent\textbf{Feature construction.}
\autoref{tab:direct_features} compares aggregation through \(Q\) with direct feature optimization from the same initialization. Both retain learned text support and matching objectives. The direct variant removes feature pooling, the residual bound, and the feature anchor. Learning rates and 30 construction updates are matched.

On SportsFit, aggregation through \(Q\) improves F1 by 5.03 points with GCN and 5.26 with GraphSAGE, with paired 95\% intervals of [4.48, 5.57] and [3.93, 6.58]. BookHis has positive mean differences, although the GCN interval includes zero. The result supports aggregation with a bounded correction under this budget.

\begin{table}[t]
\centering
\footnotesize
\setlength{\tabcolsep}{3pt}
\caption{Aggregation through $Q$ versus direct feature optimization. Macro-F1 is mean $\pm$ SD over three construction seeds on one order, each averaging three student runs. Both architectures use the same memories.}
\label{tab:direct_features}
\begin{tabular*}{\columnwidth}{@{\extracolsep{\fill}}llrr@{}}
\toprule
Dataset & Features & GCN & GraphSAGE\\
\midrule
SportsFit & Aggregation & 80.06$\pm$0.66 & 77.96$\pm$0.67\\
& Direct & 75.03$\pm$0.47 & 72.71$\pm$0.76\\
BookHis & Aggregation & 49.00$\pm$1.01 & 49.45$\pm$1.36\\
& Direct & 46.33$\pm$0.64 & 47.95$\pm$2.43\\
\bottomrule
\end{tabular*}
\end{table}

\noindent\textbf{Shared support.}
\autoref{tab:assignment_cost} tests one support matrix for both features and text with matched initialization, objectives, and budgets. Sharing halves the assignment scores and reduces construction peaks from 15.99 to 12.08 GiB on ArXiv, from 32.60 to 24.72 GiB on SportsFit, and from 3.26 to 2.77 GiB on BookHis. The corresponding construction-time reductions are 17.5\%, 21.5\%, and 18.1\%. F1 changes by +0.04, -0.37, and -1.92 points. The result supports a smaller assignment parameterization and lower construction cost, not accuracy equivalence.

\begin{table}[t]
\centering
\footnotesize
\setlength{\tabcolsep}{3pt}
\caption{Shared and separate assignments in the controlled dense implementation. F1 is mean $\pm$ SD over three construction seeds on one order. Peak allocated GPU memory and total construction time are means for those same runs.}
\label{tab:assignment_cost}
\begin{tabular*}{\columnwidth}{@{\extracolsep{\fill}}llrrr@{}}
\toprule
Dataset & Assignment & F1 & GiB & s\\
\midrule
ArXiv & Separate & 53.20$\pm$0.12 & 15.99 & 292.5\\
& Shared & 53.24$\pm$0.09 & 12.08 & 241.3\\
SportsFit & Separate & 80.70$\pm$0.66 & 32.60 & 226.7\\
& Shared & 80.33$\pm$0.34 & 24.72 & 177.9\\
BookHis & Separate & 48.47$\pm$1.18 & 3.26 & 81.6\\
& Shared & 46.55$\pm$0.93 & 2.77 & 66.9\\
\bottomrule
\end{tabular*}
\end{table}

\subsection{Memory Reuse (RQ3)}
\label{sec:student_transfer}

The memories in~\autoref{tab:direct_features} train both GCN and GraphSAGE with the same encoder, prompts, and budget. GraphSAGE obtains 77.96 F1 on SportsFit and 49.45 on BookHis. The SportsFit advantage of aggregation persists across architectures, so the memory can be reused without reconstruction.

The \emph{full-history reference} trains the same student on all observed construction nodes with the same encoder and budget. The \(5\%\) memories retain 98.23\% of reference F1 on SportsFit and 103.33\% on BookHis (\autoref{tab:retention}). A value above 100\% means that the smaller set performs better under this finite budget, not that it contains more source information.

\begin{table}[t]
\centering
\caption{Full-history utility on one fixed order. Each entry reports Macro-F1 and retention.}
\label{tab:retention}
\resizebox{\columnwidth}{!}{%
\begin{tabular}{ccccc}
\toprule
\multirow{2}{*}{Dataset} & Full & 5\% Memory & \multirow{2}{*}{F1 Retention} & \multirow{2}{*}{Accuracy Retention} \\
 & Macro-F1 & Macro-F1 & & \\
\midrule
SportsFit & 81.46 & 80.02 & 98.23\% & 95.58\% \\
BookHis & 46.96 & 48.53 & 103.33\% & 99.14\% \\
\bottomrule
\end{tabular}%
}
\end{table}

\subsection{Resource Use (RQ4)}
\label{sec:update_sensitivity}

\noindent\textbf{Construction.}
Memory construction stores assignment parameters, differentiable normalization states, and proxy intermediates. An optional implementation uses class blocks and offloads evaluation caches, reducing the SportsFit peak from 23.80 to 17.51 GiB while changing mean F1 from 79.68 to 79.41.

\noindent\textbf{Downstream training.}
The principal resource benefit is reuse of the small training set. Under the
same 2,048-update budget, memory reduces peak GPU allocation by 89.5--93.8\%
and training time by 50.9--88.4\% on ArXiv, SportsFit, and BookHis
(\autoref{tab:student_resource_audit}). Measurements use precomputed text
embeddings with the encoder off the GPU. Construction and inference have
separate costs. The detailed measurements are retained in the compact
appendix.
